% F079 有理根定理与试根分解法：简要说明（供检索与快速引用）
\textbf{结论名称：} 有理根定理与试根分解法（Rational Root Theorem and Trial-Root Factorization）

\textbf{适用条件：} $f(x)=a_nx^n+\cdots+a_1x+a_0$ 为整数系数多项式，$n\ge 1$，$a_n\ne 0$，先考虑 $a_0\ne 0$。

\textbf{核心结论：} 若最简分数 $\dfrac pq$（$p\in\mathbb Z$，$q\in\mathbb Z_{>0}$，$p,q$ 互质）是 $f(x)=0$ 的有理根，则
\[
\boxed{p\mid a_0,\qquad q\mid a_n.}
\]
即分子 $p$ 必须是常数项的约数，分母 $q$ 必须是最高次项系数的正约数。这是\textbf{必要条件}，不是充分条件。

\textbf{试根分解：} 先列 $p/q$ 的候选值并逐个代入；若 $f(p/q)=0$，则 $qx-p$ 是 $f(x)$ 的因式，且可以写成 $f(x)=(qx-p)g(x)$，其中 $g(x)$ 仍为整数系数多项式。再作多项式除法，继续分解。

\textbf{特别情况：} 若 $a_0=0$，先提出因式 $x$，记下根 $x=0$，再对剩余多项式使用上述方法。

\textbf{例：} $x^4-x^3-3x^2+x+2$ 的有理根候选为 $\pm1,\pm2$。因 $f(1)=0$，除以 $x-1$ 得 $x^3-3x-2$，继续分解可得
\[
\boxed{x^4-x^3-3x^2+x+2=(x-1)(x-2)(x+1)^2.}
\]
